\documentclass[a4paper, 12pt]{article}
\usepackage{lmodern}
\usepackage[T1]{fontenc}
\usepackage[utf8]{inputenc}
\usepackage{url}
\usepackage{xcolor}
\usepackage[square,numbers,sort&compress]{natbib}
\bibliographystyle{elsarticle-num}

\definecolor{munsell}{rgb}{0.0, 0.5, 0.69}
\newcommand\cmnt[1]{\textcolor{munsell}{#1}}

\topmargin=0cm
\topskip=0cm
\textheight=25cm
\headheight=0cm
\headsep=0cm
\oddsidemargin=0cm
\evensidemargin=0cm
\textwidth=16cm
\parindent=0cm
\parskip=12pt

\renewcommand{\baselinestretch}{1.2}

\begin{document}

%%%%%%%%%%%%%%%%%%%%%%%%%% Filled out by the candidate! %%%%%%%%%%%%%%%%%%%%%%%%%%
\newcommand{\ImeKandidata}{Luka} % Name
\newcommand{\PriimekKandidata}{Pavićević} % Surname
\newcommand{\VpisnaStevilka}{63220489   } % Enrollment number
\newcommand{\StudyProgramme}{Computer and information science, MAG} % Study programme
\newcommand{\NaslovBivalisca}{ Moškričeva 30, 1000 Ljubljana, Slovenia } % the candidate address
\newcommand{\SLONaslov}{Analiza pristopov za pridobivanje informacij iz strulturianih podatkov z uporabo
velikih jezikovnih modelov} % Slovenian title
\newcommand{\ENGNaslov}{Analysis of approaches for extracting information from structured data using large
language models} % English title
%%%%%%%%%%%%%%%%%%%%%%%%%% End of filling in  %%%%%%%%%%%%%%%%%%%%%%%%%%

\newcommand{\ProblemParagraph}{\noindent\textbf{Problem \& State of the Art.} }
\newcommand{\ContributionsParagraph}{\noindent\textbf{Expected Contributions / Technical outcome. } }
\newcommand{\MethodologyParagraph}{\noindent\textbf{Methodology \& Validation. } }

\newcommand{\Kandidat}{\ImeKandidata~\PriimekKandidata}
\noindent
\Kandidat\\

\NaslovBivalisca \\
Study programme: \StudyProgramme \\
Enrollment number: \VpisnaStevilka
\bigskip

{\bf Committee for Student Affairs}\\
Univerza v Ljubljani, Fakulteta za računalništvo in informatiko\\
Večna pot 113, 1000 Ljubljana\\

{\Large\bf
  {\centering
    The master’s thesis topic proposal \\%[2mm]
\large Candidate: \Kandidat \\[10mm]}}

I, \Kandidat, a student of the 2nd cycle study programme at the Faculty of computer and information science,
am submitting a thesis topic proposal to be considered by the Committee for Student Affairs with the following title:

%\hfill\begin{minipage}{\dimexpr\textwidth-2cm}
Slovenian: {\bf \SLONaslov}\\
English: {\bf \ENGNaslov}
%\end{minipage}

This topic was already approved last year: \textit{NO}

I declare that the mentors listed below have approved the submission of the thesis topic proposal described
in the remainder of this document.

I would like to write the thesis in English with the following reason: Because Slovenian is not my native
language, and I can express myself better in English.

I propose the following mentor:

%%%%%%%%%%%%%%%%%%%%%%%%%% Filled in by the candidate! %%%%%%%%%%%%%%%%%%%%%%%%%%
\hfill
\begin{minipage}{\dimexpr\textwidth-2cm}
  Name, surname and title: prof. dr. Damjan Vavpotič\\
  Institution: University of Ljubljana, Faculty of Computer Science \\
  E-mail: damjan.vavpotic@fri.uni-lj.si
\end{minipage}

I propose the following co-mentor:

\hfill
\begin{minipage}{\dimexpr\textwidth-2cm}
  Name, surname and title: mag. Dragoslav Radin\\
  Institution: DevRev \\
  E-mail: dr@devrev.ai
\end{minipage}
%%%%%%%%%%%%%%%%%%%%%%%%%% End of filling in %%%%%%%%%%%%%%%%%%%%%%%%%%

\bigskip

\hfill Ljubljana, \today. \newpage

\section*{Key-words}
Large Language Model, Data Serialization, Structured data, Information Extraction, Text-to-SQL

\section*{Detailed thesis proposal}

\ProblemParagraph
When working with Large Language Models, for structured data processing, we have to consider three aspects,
accuracy, latency, and cost. When providing such data in LLM's context we have to choose appropriate format,
such as JSON, YAML, TOML. While there are works that consider the accuracy, based on different data formats
\cite{struct-text-eval}, there are non that take other two aspects into consideration. Because these formats
are easily convertible between each other, by choosing the optimal one we can directly influence our cost and
latency with minimal penalty. Furthermore there are two possibilities for extracting information with LLMs,
feeding raw data directly into LLM \cite{struct-text-eval} and generating code/SQL prompts \cite{gao2023text,
li2023can, gao2023pal}. We want to see if it is more cost efficient to utilize the second method compared to
direct inference from context.

\ContributionsParagraph
We plan to propose a decision framework defining which combination of format and method yields the lowest
latency and cost for datasets of varying sizes (e.g. <1k tokens vs. >100k tokens). These results will be
validated through reproducible experiments measuring Time-to-First-Token (TTFT), total generation time, and
extraction accuracy.

\MethodologyParagraph
We will utilize a diverse dataset which will be serialized in different data formats (Markdown , JSON ,
\LaTeX , CSV , XML , YAML , TOML , TOON \cite{toon2025spec} ). We will be testing two distinct methods,
injecting serialized data into the prompt and instructing the LLM to extract answers directly. The second
method would include providing the LLM only with the limited data examples and prompting it to generate an
SQL prompt. We will measure performance across multiple LLMs (e.g. GPT-oss \cite{agarwal2025gpt}, Gemma-3
\cite{team2025gemma}, and Deepseek-R1 \cite{guo2025deepseek}) using metrics such as input/output token count,
Time-to-First-Token, and overall accuracy.


\bibliography{bib/references}

\end{document}
